\section{Introduction}
\label{sec:intro}

Probabilistic programs make random choices, so repeated runs on the same
input may behave differently. A fundamental correctness property is
\emph{almost-sure termination} (AST): a program terminates with probability
one on a given input~\cite{SahebDjahromi1978}. Infinite executions are permitted, provided
that their total probability is zero. For example, tossing a fair coin
until the first heads is almost-surely terminating, although the infinite
sequence of tails never terminates. Proving AST is challenging: deciding
AST on a fixed input is as hard as the \emph{universal halting problem},
which asks whether a deterministic program terminates on every
input~\cite{KaminskiKM2019}.
An important application is the verification of \emph{samplers}, which
generate values according to a prescribed probability distribution and
underpin simulation and probabilistic programming. Their verification
requires both the intended output distribution and termination with
probability one.

\paragraph{From control-flow graphs to loop iterations.}
We develop an AST rule within a framework for reasoning about probability
distributions over program states~\cite{zilken2026verifyingsamplingalgorithmsdistributional},
so that output correctness and termination can be established together.
McIver et al.~\cite{10.1145/3158121} give a source-level AST rule using one
nonnegative state function. It satisfies a \emph{supermartingale condition},
which prevents an increase in expectation, and a \emph{progress condition},
which ensures a positive probability of decrease.

Majumdar and Sathiyanarayana~\cite{DBLP:journals/pacmpl/MajumdarS25}
separate these roles into two functions, called certificates, and prove
soundness and relative completeness for discrete probabilistic programs
with bounded nondeterminism. Here, bounded nondeterminism permits only
finitely many alternatives at each step, and relative completeness means
that every AST program has a proof assuming access to all true arithmetic
assertions. They also establish relative completeness of McIver et al.'s
earlier rule, resolving its previously open completeness question.

Their two-certificate rule is formulated over a probabilistic control-flow
graph (pCFG), whose nodes represent program locations and whose edges
represent execution steps. Applying it requires an invariant and
certificate values at every location, including the assignments, tests,
and choices inside a loop body. The certificates therefore depend on
both the program location and the values of the variables.

Our rule instead uses the probability distribution produced by one
complete execution of the body. This combines its assignments, tests,
and choices into one semantic step. Both certificates are specified only
where the loop guard is evaluated; no intermediate annotations are needed.
We establish soundness and relative completeness for loops with loop-free
bodies and fixed choice probabilities by translating certificates between
these two levels.

\paragraph{Two certificates for almost-sure termination.}
A \emph{variant} assigns a nonnegative integer to each state. Whenever the
guard holds, one iteration must decrease it with probability at least a
fixed $\varepsilon>0$. Other outcomes may increase it, for example after
rejection, so progress alone does not guarantee AST.
A second, nonnegative \emph{supermartingale} function must not increase in
expectation. We link the functions by requiring the variant to be bounded
on every set where the supermartingale is bounded by a fixed constant;
the variant's bound may depend on that constant. Both functions are
positive while the guard holds and zero when it does not.

Intuitively, the supermartingale condition ensures that almost every
execution stays within some bounded range of its values. The coupling
then bounds the variant along that execution. Within a fixed bound,
sufficiently many consecutive decreases lead to termination, and the
uniform progress probability prevents the variant from staying positive
forever.
A distributional invariant describes possible distributions at loop heads.
It supplies the states on which these conditions must hold: those assigned positive mass by one of its distributions.
Its quantitative information remains available for output correctness.

\paragraph{Illustration: the Fast Dice Roller.}
Lumbroso's Fast Dice Roller
(FDR)~\cite{DBLP:journals/corr/abs-1304-1916} simulates a fair $n$-sided
die using independent fair bits, returning each integer in
$\{0,\ldots,n-1\}$ with equal probability.
It maintains a candidate $c$, initially zero, and a range size $v$,
initially one. The candidate is uniform on $\{0,\ldots,v-1\}$.
Each iteration draws a fresh bit $b\in\{0,1\}$ and updates
$c\gets2c+b$ and $v\gets2v$: appending a bit doubles the candidate range.
If $v<n$, another bit is needed. Once $v\geq n$, the algorithm returns
$c$ if $c<n$; otherwise it subtracts $n$ from both $c$ and $v$ and
continues. This \emph{recycling} retains randomness from rejected candidates.

For $n=6$, the first three bits produce eight equally likely candidates.
The values $0,\ldots,5$ are accepted. Conditional on rejection, $c$ is
uniform on $\{6,7\}$; subtracting six leaves a uniform candidate in
$\{0,1\}$ with range size two. This retains one fair bit, so only two
additional flips are needed to obtain eight equally likely candidates again.

Our termination proof treats range expansion, the random choice, and
possible recycling as one iteration. Its successor distribution must
still be derived from the body, but certificate values are compared only
at loop heads. The explicit pCFG proof in \Cref{app:fdr-pcfg-proof}
shows the additional annotations at intermediate locations.

\paragraph{Contributions.}
The main contributions of this paper are:
\begin{enumerate}
\item \textbf{An AST proof rule using distribution transformers.}
We formulate a two-certificate rule for probabilistic loops with loop-free
bodies (\Cref{sec:rule}) and prove soundness and relative completeness
by translating certificates to and from the pCFG rule of Majumdar and
Sathiyanarayana~\cite{DBLP:journals/pacmpl/MajumdarS25}
(\Cref{sec:metatheory}).

\item \textbf{AST proofs for discrete sampling algorithms.}
We apply the rule to FDR, the Fast Loaded Dice Roller (FLDR), and
Han--Hoshi sampling
(\Cref{sec:case-studies}). Prior work already establishes AST of FDR
and FLDR using pCFG certificates alongside distributional
invariants~\cite{zilken2026verifyingsamplingalgorithmsdistributional}.
We revisit these arguments through the general iteration-level rule
and make their support assertions explicit. For FDR, we compare the
same logarithmic progress argument at both levels; a simpler affine
variant is also available at either level. Han--Hoshi adds an application
based on interval refinement.

\item \textbf{Template-based certificate synthesis.}
Our prototype searches guarded affine and piecewise affine templates
using supplied invariants and iteration semantics
(\Cref{sec:certificate-synthesis}). Bounded integer coefficients and
fixed partitions make the search incomplete for AST programs.
Experiments on an unbounded random walk, symbolic FDR, a concrete FLDR
tree, and a structural FLDR abstraction demonstrate synthesis of both
certificates, including for a loop with infinite expected runtime.
\end{enumerate}

\paragraph{Scope.}
We consider fixed-probability discrete choices without nondeterminism
and unnested while-loops, i.e., loops with loop-free bodies. Bounded body
path length and fixed probabilities give a uniform positive lower bound
on every feasible body path. The sampling loops have this form, with
terminating preprocessing treated separately. Nested loops and
nondeterminism are beyond the present metatheory. Relative completeness
concerns the existence of a suitable invariant and certificates; it does
not imply completeness of the restricted template search.

\paragraph{Organisation of the Paper.}
\Cref{sec:prelims} introduces the semantics and distributional invariants.
\Cref{sec:rule} presents our rule with random-walk and FDR illustrations.
\Cref{sec:metatheory} recalls the complete pCFG rule, compares the two
interfaces, and proves soundness and relative completeness.
\Cref{sec:case-studies} develops the sampler proofs;
\Cref{sec:certificate-synthesis} presents synthesis and its results.
\Cref{sec:relwork,sec:conclusion} discuss related work and conclude.
